This report presents a versioned study of an AI-based phishing detector. The first experiment independently reproduced a character-level CNN and 36-feature LightGBM ensemble on 20,000 URLs. The deployed application later adopted a character CNN, two-layer bidirectional LSTM and additive attention. A source-level audit found that the newer trainer did not preserve the original root-domain separation guarantee. Whole-domain reservations, consistent cleaned-input identity and frozen split hashes were therefore introduced before a separately versioned v3.4 training run.

The corrected experiment retains 63,869 records: 25,961 training, 6,113 validation and 31,795 test records. All pairwise hostname and registered-domain intersections are zero. At threshold 0.5, the candidate achieves 92.876\% test accuracy, 90.696\% recall and 98.289\% ROC-AUC, with 718 false positives and 1547 false negatives. The frozen validation gate failed. The current deployed model remains v3.3; candidate research results are not presented as deployed performance.

Independent HTML and email analyzers provide contextual warning signals without visiting destinations, executing pasted HTML or accessing a mailbox. Human-verified feedback, model versioning, deterministic robustness tests and drift-monitoring infrastructure extend the initial URL-only reproduction. Their scope and limitations are documented separately from classifier accuracy. The report emphasizes reproducible evidence and honest trade-offs rather than assuming that more data or a newer architecture necessarily improves every metric.